
\label{app:data-sources}

All datasets and reference parameters used in this thesis are publicly
accessible. The primary experimental dataset, and the contextual
reference sources used to motivate and parameterise the simulation, are
listed below.

\begin{itemize}[leftmargin=*]
    \item \textbf{Primary experimental dataset}:
    ITU 2024 Smart Energy Scheduling Challenge dataset (ten telecom base
    station sites, 60 days each, hourly resolution), released via Zindi
    \citep{itu2024zindi} —
    \url{https://zindi.africa/competitions/smart-energy-supply-scheduling-for-green-telecom}

    \item \textbf{Solar irradiance and temperature}:  
    NREL NSRDB PSM~v3 — \url{https://nsrdb.nrel.gov}

    \item \textbf{Telecom load references}:  
    TRAI QoS and performance indicator reports — \url{https://www.trai.gov.in}

    \item \textbf{Grid reliability statistics}:  
    Central Electricity Authority (CEA) annual reports — \url{https://cea.nic.in} \\
    UDAY DISCOM portal — \url{https://www.uday.gov.in}

    \item \textbf{Diesel logistics}:  
    GSMA Green Power for Mobile reports — \url{https://www.gsma.com} \\
    DoT rural audits and energy reports — \url{https://dot.gov.in}

    \item \textbf{Battery parameters}:  
    LFP manufacturer datasheets (Delta, Huawei, Emerson, Amara Raja)

    \item \textbf{Economic and emission factors}:  
    IOCL/BPCL/HPCL fuel-price bulletins — \url{https://iocl.com} \\
    State Electricity Regulatory Commissions (SERC) tariff orders \\
    CEA CO$_2$ Baseline Database — \url{https://cea.nic.in}
\end{itemize}
